%!TEX root = ../main.tex

\begin{abstract}

Question answering over documents in a dialogue is harder than answering single
questions: the answers have to come from the documents, follow-up questions depend on
earlier turns, and some questions cannot be answered from the documents at all. This
work compares several approaches to this task on \gls{mdd}, a dataset of goal-oriented
dialogues about US government services in four domains.

The work starts from an \gls{llm} used alone and adds retrieval step by step. The
\gls{llm} alone reaches 15.9\% token \fone{} on the validation split. Retrieval-augmented
generation with hybrid sparse and dense retrieval and an \gls{llm} rewriting of the
follow-up questions reaches 26.8\%, and query rewriting is the only change with a clearly
significant effect ($+2.0$ points over hybrid retrieval, $p < 0.001$). When the model is
given the correct sections directly, it reaches 36.1\%, and about two thirds of the gap
to this oracle comes from retrieval errors. A cross-encoder reranker raises \ratk{5} from
61.2\% to 67.4\% on a 500-question development subset, but the answers do not improve
significantly: the reranker finds new gold sections but also pushes others out of the
top five.

A knowledge graph is then built from the corpus with an \gls{llm}, with three layers for
the document structure, the extracted mentions and the resolved entities. The extraction
is noisy and not deterministic (two runs with the same settings share fewer than half of
their triples), and embedding similarity alone is not enough to decide which entities to
merge. On a small annotated sample a single extraction run has 51\% precision and 66\%
recall. The final graph has 14{,}731 entities and 13{,}277 relations, covers the gold
section of 96.6\% of the development questions, and has a precision of 73\% on a control
sample.

For retrieval, the rewritten user turn is linked to graph entities, which reaches an
entity of the gold section for 70\% of the questions. Retrieval over the graph alone,
however, finds the gold section only half as often as passage retrieval, mainly because
entity overlap cannot rank the many sections the linked entities lead to. Combining
graph and passage retrieval matches the passage pipeline without improving it, while
adding graph facts to the retrieved passages gives a small but consistent gain. The
remaining choices were settled on the full validation split, and the final systems were
run once on the test split: the \gls{llm} alone reaches 15.1\% \fone{}, graph retrieval
23.8\%, the best passage pipeline 26.9\%, passages with graph facts 27.5\% and the oracle
34.2\%. A control with the same number of random facts shows that the gain of 0.8 points
on answerable questions comes from the content of the facts and not from the longer
context. Giving the generator the dialogue history helps on answerable questions but makes
it abstain less; with the decision to abstain taken without the history, passages with
graph facts reach 28.0\% on validation. Two newer generator models answer as well when they
answer, but refuse more answerable questions, so \texttt{gpt-5.4-nano} was kept. The graph facts help mostly where passage retrieval misses
the answer and a fact still comes from the right section, and a router using signals
available before generation could not tell these questions apart.

\end{abstract}
